% ECONOMICS_PROBLEM: {"id": "C71-299", "title": "A combinatorial polynomial-time nucleolus algorithm for matching games", "jel_code": "C71", "source_id": "OP-0201"}
% !TeX program = xelatex
\documentclass[11pt,a4paper]{article}
\usepackage[margin=23mm]{geometry}
\usepackage{fontspec}
\setmainfont{Latin Modern Roman}
\usepackage{amsmath,amssymb,mathtools}
\usepackage{xcolor,enumitem,booktabs,longtable,array,xurl,hyperref}
\definecolor{muted}{HTML}{53636C}
\hypersetup{colorlinks=true,urlcolor=blue,linkcolor=blue}
\setlength{\parindent}{0pt}
\setlength{\parskip}{5pt}
\newcommand{\R}{\mathbb R}
\newcommand{\E}{\mathbb E}
\newcommand{\N}{\mathbb N}
\newcommand{\ind}{\mathbf 1}
\newcommand{\Var}{\operatorname{Var}}
\newcommand{\Cov}{\operatorname{Cov}}
\newcommand{\supp}{\operatorname{supp}}
\DeclareMathOperator*{\argmin}{arg\,min}
\DeclareMathOperator*{\argmax}{arg\,max}
\newcommand{\entrymeta}[3]{{\sffamily\small\color{muted}\textbf{\texttt{#1}}\quad|\quad #2\par #3\par}}
\newcommand{\Problemref}[1]{\texttt{#1}}
\newcommand{\JELref}[2]{\texttt{#2} (JEL \texttt{#1})}
\begin{document}
\section*{A combinatorial polynomial-time nucleolus algorithm for matching games}
\textbf{Problem ID:} \texttt{C71-299}\quad \textbf{JEL:} \texttt{C71}\par
% BEGIN ECONOMICS STATEMENT
\label{problem:OP-0201}
% GAME_THEORY_PROBLEM: {"local_id": "GT-071", "original_number": 71, "kind": "P", "entry_kind": "statement_only", "published": false, "review_required": true, "category": "game_theory"}
\label{game:GT-071}
{\sffamily\small\color{muted}
\noindent\textbf{ID: GT-071.}\quad\textbf{Category:} Cooperative game algorithms.
\quad\textbf{Kind: P.}\quad\textbf{Status:} Open algorithm-design question in the cited source.
\par}
\subsection*{Definitions and mathematical statement}
Let \(G=(V,E)\) be a finite simple undirected graph, with rational
weights \(w_e\geq 0\). The players are the vertices. For \(S\subseteq V\),
define the worth of coalition \(S\) by
\[
 \nu(S)=
 \max\left\{\sum_{e\in F}w_e:
              F\text{ is a matching in the induced graph }G[S]\right\},
\]
where the empty matching is permitted. An imputation is
\[
 x\in\mathcal I(\nu)
   :=\left\{x\in\mathbb R_{\geq 0}^{V}:
                         \sum_{i\in V}x_i=\nu(V)\right\}.
\]
For each nonempty proper coalition \(S\), its excess at \(x\) is
\(e(S,x)=\nu(S)-\sum_{i\in S}x_i\).
Let \(E(x)\) be the vector of all these excesses, sorted in
nonincreasing order, with multiplicities. The nucleolus is the unique
imputation \(x^{\mathrm{nuc}}\) that lexicographically minimizes
\(E(x)\) over \(\mathcal I(\nu)\).

\paragraph{Question.}
Is there a deterministic algorithm computing \(x^{\mathrm{nuc}}\)
in time polynomial in the binary input length using explicit
combinatorial graph operations and rational arithmetic, without a
general linear-programming solver or the ellipsoid method with
black-box separation?

\subsection*{Short English statement}
Can the nucleolus of a general weighted matching game be computed by
a polynomial-time combinatorial algorithm?

\subsection*{Variants and scope boundaries}
The graph need not be bipartite. Polynomial-time computability is already
known; the restriction concerns the method of computation and is not
a standard complexity-class separation. Strongly polynomial time is a
stronger optional objective. The core can be empty, so the nucleolus
must be defined on imputations, rather than on the core alone.

\subsection*{Source relationship and status}
[S1, Section 2] states this algorithm-design problem.
Its definition of a combinatorial algorithm excludes an LP solver.
The lexicographic objective above is over coalition excesses, not
over individual agents' allocations.

\subsection*{References}
\begingroup\footnotesize\raggedright
\noindent[S1] Vijay V.~Vazirani,
\emph{Equitable Core Imputations via a New Adaptation of the
Primal-Dual Framework}, 2024, arXiv:2402.11437v4.
\url{https://arxiv.org/html/2402.11437v4}.
Locators: Section 1, footnote 3 (the algorithmic convention);
Section 2, paragraphs following Definition 1 (nucleolus computation
and the open combinatorial problem).
\par\endgroup

\subsection*{Common conventions: Source mathematical conventions}

\textbf{Finite games.} For a finite set $A$, $\Delta(A)$ is its probability
simplex. Players choose independent mixed strategies in Nash equilibrium;
correlation is allowed only when explicitly part of the solution concept.
In a game with payoff functions $u_i$ and strategy profile $x$, an additive
$\varepsilon$ Nash equilibrium satisfies
\[
 u_i(x)\geq u_i(x_i',x_{-i})-\varepsilon
 \quad\text{for every player $i$ and every admissible deviation $x_i'$}.
\]
For numerical approximation, payoffs are normalized as locally specified.
An $\varepsilon$ payoff residual is distinct from distance at most
$\varepsilon$ to the set of exact equilibria. Point convergence, convergence
of empirical play, and convergence in distance to an equilibrium set are
also distinct.

\textbf{Computational representations.} Unless stated otherwise, a complexity
question takes rational data with binary encoding; polynomial time is measured
in total bit length. A fixed-accuracy polynomial algorithm is not necessarily
polynomial in $1/\varepsilon$, and neither is necessarily polynomial in
$\log(1/\varepsilon)$. Succinct games and value, demand, or sampling oracles
must declare their interfaces. Exact real solutions require an explicit
representation; an unspecified list of arbitrary real numbers is not a
finite computational output. A claimed algorithmic barrier is meaningful
only for its specified model and reductions.

\textbf{Dynamic games.} Histories, information, observation, and allowable
randomization are part of the model. A stationary strategy depends only on
the current state; a positional strategy in a deterministic graph game
chooses one action at each controlled vertex. Nash equilibrium at an initial
state and equilibrium after every history are different requirements.
An expectation of a pathwise $\liminf$ payoff, a $\liminf$ of expected averages,
a limiting expected average, and a uniform finite-horizon equilibrium are
different criteria. None is silently substituted for another.

\textbf{Allocation and mechanisms.} A complete allocation assigns every item
exactly once unless otherwise stated. Zero-valued goods, negative values,
capacity constraints, feasibility, lotteries, and copies of goods affect
fairness definitions and are specified locally. Dominant-strategy incentive
compatibility, Bayesian incentive compatibility, universal truthfulness,
and truthfulness in expectation impose different inequalities. Whenever
an objective ratio has zero denominator, its convention is stated or those
instances are excluded.

\textbf{Infinite-dimensional models.} Expectations, integrals, stochastic
equations, and optimization objectives require their local regularity,
integrability, filtrations, and admissible domains. A backward--forward
mean-field system is a boundary-value problem; it is not an initial-value
semiflow without an additional construction. Long-time, large-population,
and vanishing-noise limits require an order of limits and a convergence
topology. If a minimum need not be attained, the objective is its infimum
or supremum and the approximation requirement is stated separately.
% END ECONOMICS STATEMENT
\end{document}
